\section*{Chapitre \ref{ch:g9:climate-past-and-present} --- Le climat de la Terre, hier et aujourd'hui}

\begin{solution}{exo:g9:climate-past-and-present:1}
(a) Environ \qty{-19}{\degreeCelsius}. (b) Environ \qty{15}{\degreeCelsius}.
(c) $15 - (-19) = 34$~: environ \qty{34}{\degreeCelsius}. (d) Deux gaz
parmi~: le dioxyde de carbone, la vapeur d'eau, le méthane.
\end{solution}

\begin{solution}{exo:g9:climate-past-and-present:2}
(a) Temps. (b) \omterm{def:g5:weather-and-climate:climate}{Climat}. (c) \omterm{def:g5:weather-and-climate:climate}{Climat} (et \omterm{def:g9:climate-past-and-present:climate-change}{changement climatique}~: une
modification sur des décennies). (d) Temps.
\end{solution}

\begin{solution}{exo:g9:climate-past-and-present:3}
(a) \qty{19}{cm}, onze années. (b) $19 \div 11 \approx \qty{1.7}{cm}$.
(c) La \omterm{def:g4:clouds-rain-snow:precipitation}{neige} de chaque année tombe sur celle des années précédentes, si
bien que plus une couche est profonde, plus elle est tombée tôt.
\end{solution}

\begin{solution}{exo:g9:climate-past-and-present:4}
(a) $22.4 - 21.0 = +\qty{1.4}{\degreeCelsius}$. (b) $20.3 - 21.0 =
\qty{-0.7}{\degreeCelsius}$. (c) \qty{0}{\degreeCelsius}. (d) Ce mois a
été plus froid que sa moyenne de 1991--2020.
\end{solution}

\begin{solution}{exo:g9:climate-past-and-present:5}
(a) Des années 1880 aux années 1930, les années 1950 et les années 1960.
(b) Environ \qty{0.25}{\degreeCelsius} et \qty{0.80}{\degreeCelsius}. (c) $0.80 -
0.25 = \qty{0.55}{\degreeCelsius}$. (d) Non~: des années 1880 aux années
1910, les décennies se sont un peu refroidies, et les années 1950 ont été
un peu plus froides que les années 1940~; mais depuis les années 1970,
chaque décennie a été plus chaude que la précédente.
\end{solution}

\begin{solution}{exo:g9:climate-past-and-present:6}
(a) $(324.62 - 315.98) \div 10 \approx \qty{0.86}{ppm}$ par an. (b)
$(427.35 - 401.01) \div 10 \approx \qty{2.63}{ppm}$ par an. (c) $2.63
\div 0.86 \approx 3.1$~: environ trois fois plus vite. (d) $(427.35 - 315.98)
\div 315.98 \times 100 \approx 35~\%$.
\end{solution}

\begin{solution}{exo:g9:climate-past-and-present:7}
(a) $8 \times 10^5$ ans. (b) Il y a environ \num{120000} à \num{130000}
ans~: la courbe des températures y passe au-dessus de 0. (c) Environ
\qty{275}{ppm}. (d) Oui~: les pics et les creux des deux courbes arrivent
aux mêmes moments.
\end{solution}

\begin{solution}{exo:g9:climate-past-and-present:8}
(a) $3.2 \times 10 = \qty{32}{mm} = \qty{3.2}{cm}$. (b) $2100 - 2025 =
75$ ans~; $75 \times 3.2 = \qty{240}{mm}$, environ \qty{24}{cm} (et
davantage si la hausse s'accélère). (c) L'eau de \omterm{def:g2:water-on-earth:ocean}{mer} se dilate en se
réchauffant~; la glace des terres fond et son eau coule dans la \omterm{def:g2:water-on-earth:ocean}{mer}.
(d) Non~: la \omterm{def:g8:ice-on-earth:sea-ice}{banquise} flotte déjà sur la \omterm{def:g2:water-on-earth:ocean}{mer}, donc sa fonte ne la fait pas
monter.
\end{solution}

\begin{solution}{exo:g9:climate-past-and-present:9}
(a) 90. (b) 66. (c) 99. (d) \emph{Probable} signifie au moins 66 chances
sur 100, pas une certitude~: le journal rend l'affirmation plus sûre que
ne l'ont dit les scientifiques.
\end{solution}

\begin{solution}{exo:g9:climate-past-and-present:10}
Un hiver froid en un seul lieu relève du \emph{temps}~: il dit ce qui
s'est passé pendant quelques mois en un seul endroit. Le \omterm{def:g9:climate-past-and-present:global-warming}{réchauffement
climatique} est une modification du \emph{\omterm{def:g5:weather-and-climate:climate}{climat}}, mesurée sur des
décennies, et \emph{planétaire}~: c'est la moyenne sur la Terre entière.
Un hiver froid ici peut se produire alors que la Terre dans son ensemble
est plus chaude qu'avant, et les moyennes par décennie continuent de
monter.
\end{solution}

\begin{solution}{exo:g9:climate-past-and-present:11}
(a) $280 - 190 = \qty{90}{ppm}$, en environ $2 \times 10^4$ ans. (b) $427
- 280 = \qty{147}{ppm}$, en $2025 - 1750 = 275$ ans. (c) La seconde
hausse est plus grande (147 contre 90~ppm). Rythmes~: $90 \div (2 \times 10^4) =
4.5 \times 10^{-3}$~ppm par an, contre $147 \div 275 \approx 5.3 \times
10^{-1}$~ppm par an~: environ 120 fois plus vite. (d) Qu'en 2019, il
était plus élevé qu'à n'importe quel moment depuis au moins deux millions
d'années (degré de confiance élevé).
\end{solution}

\begin{solution}{exo:g9:climate-past-and-present:12}
(a) Ils sont presque égaux~: 1.07 contre \qty{1.06}{\degreeCelsius}. La
meilleure estimation dit que l'influence humaine explique à peu près la
totalité du réchauffement observé. (b) Non~: les causes naturelles réunies
ont modifié la température de \qty{0.1}{\degreeCelsius} au plus, dix fois
moins que le réchauffement observé. (c) Au moins 66 chances sur 100 que le
réchauffement d'origine humaine soit compris entre 0.8 et
\qty{1.3}{\degreeCelsius}.
\end{solution}

\begin{solution}{pb:g9:climate-past-and-present:1}
\textbf{1.} Parce que c'est la période de référence~: les anomalies sont
mesurées à partir de sa propre moyenne.

\textbf{2.} Plus froide, de \qty{0.22}{\degreeCelsius}.

\textbf{3.} Plus chaude, de \qty{0.98}{\degreeCelsius}.

\textbf{4.} $0.98 - (-0.22) = \qty{1.20}{\degreeCelsius}$.

\textbf{5.} $427 - 316 = \qty{111}{ppm}$.

\textbf{6.} $111 \div 316 \times 100 \approx 35~\%$.

\textbf{7.} $427 - 300 = \qty{127}{ppm}$.

\textbf{8.} Le dioxyde de carbone absorbe une partie de la lumière
infrarouge émise par la surface chaude et en renvoie une part vers le bas.
Quand l'air en contient davantage, plus d'énergie revient vers la surface,
qui se réchauffe.

\textbf{9.} $2025 - 1993 = 32$ ans~; $32 \times 3.2 \approx
\qty{102}{mm}$, environ \qty{10}{cm}.

\textbf{10.} L'eau de \omterm{def:g2:water-on-earth:ocean}{mer} qui se réchauffe se dilate~; la glace des terres
(\omterm{def:g8:ice-on-earth:glacier}{glaciers}, \omterm{def:g8:ice-on-earth:ice-sheet}{calottes glaciaires}) fond et rejoint la \omterm{def:g2:water-on-earth:ocean}{mer}.

\textbf{11.} Oui, du même ordre~: environ \qty{1.2}{\degreeCelsius} contre
\qty{1.09}{\degreeCelsius}. Ils diffèrent un peu parce que les périodes ne
sont pas les mêmes (1880--1899 au lieu de 1850--1900~; 2015--2024 au lieu
de 2011--2020~: les années les plus récentes ont été plus chaudes), et
parce que le GIEC combine plusieurs séries, pas seulement celle de la
NASA.

\textbf{12.} «~Depuis la fin du \textsc{xix}\textsuperscript{e}~siècle, la
surface de la Terre s'est réchauffée d'environ
\textbf{\qty{1.2}{\degreeCelsius}}.~»
\end{solution}
